% ===================================================================
%  圖表第 4 號 — 壯年同老年。
%
%  Traduction, faite en 2026, en cantonais ecrit d'aujourd'hui, en
%  caracteres traditionnels, sur l'ido de texto/io. Regles de la
%  colonne : voir 10-toubiu-01.tex. Deux scenes, deux numerotations :
%  les cles portent c1 et c2.
%
%  LA PARENTE CANTONAISE EST LA PLUS FINE DE TOUTE LA SERIE, et ce
%  tableau la met entiere sur une seule table. La ou le vietnamien
%  distinguait le cote et l'age, le cantonais distingue aussi le COTE
%  PAR ALLIANCE, et il le fait avec des mots differents pour le mari
%  et pour la femme :
%
%    家公 / 家婆   les beaux-parents DE L'EPOUSE
%    externe : 外父 / 外母  les beaux-parents DE L'EPOUX
%    心抱          la bru      (Pekin : 儿媳)
%    女婿          le gendre
%    大伯 / 叔仔   le frere aine / cadet du mari
%    姑奶 / 姑仔   la soeur ainee / cadette du mari
%
%  « 心抱 » est le mot le plus cantonais du tableau : Pekin ecrit
%  儿媳, et personne a Hong Kong ne le dit. Il se prononce « sam1
%  pou5 » et s'ecrit avec le coeur.
%
%  ET LE REPAS DE NOCES OBLIGE A LIRE LA PLANCHE, comme il l'avait
%  fait en vietnamien : c'est la mere DU MARIE qui sourit (17), le
%  beau-pere DE LA MARIEE qui est heureux d'avoir une 心抱 (20). Le
%  cantonais ne peut pas ecrire ces phrases sans avoir tranche.
%
%  LES MALADIES DE LA GRAND-MERE SONT CELLES QU'ON DIT A HONG KONG :
%  頭暈 le vertige, 傷風 le rhume, 牙痛 le mal de dents, 風濕 le
%  rhumatisme, 肺炎 la pneumonie. Aucune n'est un emprunt ; le
%  vocabulaire de la maladie est le plus ancien de la langue.
% ===================================================================

%%K t04-apar-1 apar
\VUsaut{4.0mm}
\VUtitre{15.0pt}{60}{圖表第 4 號}
\VUsaut{1.6mm}
\VUtitre{11.4pt}{0}{壯年同老年。}
\VUsaut{3.2mm}

%%K t04-tit-1 sub
\VUcentre{11.4pt}{0}{\textit{第一場。}}
\VUsaut{1.6mm}
\VUtitre{11.4pt}{0}{喜酒。}
\VUsaut{2.6mm}

%%K t04-tit-2 sub
\VUpk{10.2pt}{40}{秋天。}
\VUsaut{3.0mm}

%%K t04-c1-01-1 p
1. --- 啲後生仔大咗喇。其中一個，約翰，娶老婆。我哋而家喺\VUgras{喜酒}度，兩家人
坐埋同一張檯。

%%K t04-c1-02-1 p
2. --- 我哋而家係\VUgras{秋天}，喺\VUgras{九月}。\VUgras{十月}同\VUgras{十一月}嘅
冷風仲未將鄉下嘅綠葉吹得晒。夜晚嘅空氣溫暖又香；所以晚飯就擺喺花園嗰邊嘅
\VUgras{騎樓底} \textsuperscript{(8)}。

%%K t04-c1-03-1 p
3. --- 遠處嘅\VUgras{山崗} \textsuperscript{(3)} 上面，就係約翰父母嘅屋，一座掩喺
綠蔭裏面嘅雅致\VUgras{城堡} \textsuperscript{(1)}；山坡上面滿係靚葡萄園，出好靚
嘅酒。種葡萄嗰個打理住佢哋。

%%K t04-c1-04-1 p
4. --- 葡萄園上面飛過一群\VUgras{鷓鴣} \textsuperscript{(2)}，佢哋畀行近嘅
\VUgras{看守} \textsuperscript{(5)} 嚇親。呢位看守挾住支\VUgras{槍}，孭住個
\VUgras{獵袋}，搜緊啲\VUgras{矮樹叢} \textsuperscript{(4)}，佢隻\VUgras{狗}
\textsuperscript{(6)} 就喺佢身邊。佢睇住呢片地，唔畀偷獵嘅人殺絕啲野味。

%%K t04-c1-05-1 p
5. --- \VUgras{騎樓底}外面爬住一棚葡萄，掛住一串串沉甸甸嘅\VUgras{葡萄}
\textsuperscript{(7)}。

%%K t04-c1-06-1 p
6. --- 裝飾住\VUgras{張檯} \textsuperscript{(16)} 嘅靚秋花係\VUgras{菊花}
\textsuperscript{(9)}；新娘個\VUgras{老豆} \textsuperscript{(10)} 摘咗一朵插喺自己
禮服嘅鈕門度。

%%K t04-c1-07-1 p
7. --- 餐飯就嚟食完。\VUgras{侍應} \textsuperscript{(11)} 開始攞甜品出嚟，一陣間就
會收埋啲唔再用嘅餐具：\VUgras{匙羹} \textsuperscript{(14)}、\VUgras{叉}
\textsuperscript{(12)}、\VUgras{刀} \textsuperscript{(13)} 同\VUgras{大碟}
\textsuperscript{(15)}。

%%K t04-tit-3 sub
\VUpk{10.2pt}{40}{親戚。}
\VUsaut{3.0mm}

%%K t04-c1-08-1 p
8. --- 個個都好開心；新郎個\VUgras{阿媽} \textsuperscript{(17)} 望住自己啲仔女嗰個
笑容，就證明佢好幸福。

%%K t04-c1-09-1 p
9. --- 呢頭婚事將本地兩家有錢人結埋一齊。約翰，就係個\VUgras{新郎}
\textsuperscript{(18)}，係一個大地主嘅\VUgras{仔}，佢而家做咗一個能幹嘅工業家嘅
\VUgras{女婿}。

%%K t04-c1-10-1 p
10. --- 呢對夫婦好後生，不過佢哋好耐之前已經\VUgras{訂咗婚}。\VUgras{新娘}
\textsuperscript{(19)} 瑪爾達，著住鑲住橙花嘅\VUgras{白色婚紗}，好迷人。

%%K t04-c1-11-1 p
11. --- 佢個\VUgras{家公} \textsuperscript{(20)} 好開心有咁乖嘅\VUgras{心抱}，
而約翰個\VUgras{外母} \textsuperscript{(21)} 就見到自己個\VUgras{女}咁靚，笑咗。

%%K t04-c1-12-1 p
12. --- 檯尾坐住其餘嘅\VUgras{客} \textsuperscript{(22)}：\VUgras{親戚、朋友、
表兄弟、表姊妹}。一個\VUgras{管酒席嘅} \textsuperscript{(23)}，餐巾搭喺手臂上面，
好殷勤咁招呼佢哋。

%%K t04-tit-4 sub
\VUpk{10.2pt}{40}{朋友。}
\VUsaut{3.0mm}

%%K t04-c1-13-1 p
13. --- 檯右邊，就係一位\VUgras{小姐} \textsuperscript{(24)}，係新娘嘅\VUgras{朋友}；
跟住係佢阿哥，做\VUgras{伴郎} \textsuperscript{(25)}。佢同自己嘅\VUgras{伴娘}
\textsuperscript{(26)} 傾偈，佢係新郎個細妹，因為呢頭婚事而做咗瑪爾達嘅
\VUgras{姑仔}。佢係個迷人嘅少女。佢有靚\VUgras{首飾}、一串\VUgras{珍珠鏈}同一隻
金\VUgras{手鈪}。

%%K t04-c1-14-1 p
14. --- 喺佢哋附近，企起身舉起酒杯嗰位先生係屋企嘅一位\VUgras{朋友}
\textsuperscript{(27)}。佢係\VUgras{新郎嘅證婚人}之一。佢\VUgras{祝酒}，飲杯酒
祝新人\VUgras{身體健康}，又祝佢哋長久幸福。佢著住禮服：燕尾服同\VUgras{漆皮鞋}
\textsuperscript{(28)}。就係佢陪住已經\VUgras{守寡}嘅\VUgras{嫲嫲} \textsuperscript{(29)}。

%%K t04-c1-15-1 p
15. --- 著住\VUgras{燕尾服} \textsuperscript{(31)}、留住幼幼黑鬚嘅後生仔，係約翰嘅
\VUgras{姐夫} \textsuperscript{(30)}。佢係個出色嘅工程師。佢隔籬坐住一位好雅致嘅
\VUgras{少婦} \textsuperscript{(33)}，佢係瑪爾達嘅\VUgras{家姐}，亦係個嬌俏女仔嘅
阿媽——嗰個女仔而家挽住表哥隻手，行嚟送花同\VUgras{祝}佢\VUgras{晚安}。呢位太太
戴住好靚嘅\VUgras{耳環}，梳住雅致嘅\VUgras{髮髻}。

%%K t04-c1-16-1 p
16. --- 個細表哥著住件\VUgras{罩衫} \textsuperscript{(36)} 同條鬆身\VUgras{短褲}
\textsuperscript{(37)}，樣衰衰，好值得可憐。佢係\VUgras{孤兒} \textsuperscript{(35)}。
佢好細個就冇咗父母。佢阿叔做咗佢嘅\VUgras{監護人} \textsuperscript{(38)}，仲為咗養大
呢個可憐嘅細路而唔娶老婆。佢係個好人。佢有啲禿頭，又好深近視；佢戴住
\VUgras{夾鼻眼鏡}。

%%K t04-tit-5 sub
\VUsaut{2.4mm}
\VUpk{10.2pt}{40}{甜品。}
\VUsaut{3.0mm}

%%K t04-c1-17-1 p
17. --- 喺啲食客後面，一張鋪住\VUgras{檯布}嘅檯上面，\VUgras{甜品}
\textsuperscript{(39)} 已經擺好。一個大杯入面係生果：\VUgras{桃}
\textsuperscript{(40)}、\VUgras{梨} \textsuperscript{(41)}、\VUgras{蘋果}
\textsuperscript{(42)}、\VUgras{杏} \textsuperscript{(43)} 同\VUgras{葡萄}
\textsuperscript{(44)}。隔籬係\VUgras{菠蘿} \textsuperscript{(45)}、一個靚
\VUgras{蛋糕} \textsuperscript{(46)}、幾樽\VUgras{烈酒} \textsuperscript{(48)} 同
\VUgras{香檳} \textsuperscript{(49)}，仲有一疊疊乾淨嘅\VUgras{碟}
\textsuperscript{(47)}。

%%K t04-c1-18-1 p
18. --- \VUgras{侍酒} \textsuperscript{(50)} 開緊一\VUgras{樽} \textsuperscript{(52)}
陳年酒，用嘅係\VUgras{開瓶器} \textsuperscript{(51)}。個\VUgras{軟木塞} \textsuperscript{(53)} 睇嚟
好難拔。呢位侍酒手臂上面搭住一條\VUgras{餐巾} \textsuperscript{(54)}。

%%K t04-c1-19-1 p
19. --- 食完晚飯，啲男士會去吸煙室，啲太太就去打扮，準備跳舞。

%%K t04-tit-6 sub
\VUsaut{5.0mm}
\VUcentre{11.4pt}{0}{\textit{第二場。}}
\VUsaut{1.6mm}
\VUpk{11.4pt}{40}{阿爺嘅生日。}
\VUsaut{2.6mm}

%%K t04-tit-7 sub
\VUpk{10.2pt}{40}{探訪。}
\VUsaut{3.0mm}

%%K t04-c2-01-1 p
1. --- 十年過去咗。約翰嘅父母老咗，好錫自己三個孫，兩個\VUgras{孫仔}
\textsuperscript{(2)} 同一個\VUgras{孫女} \textsuperscript{(3)}。因為今日係
\VUgras{阿爺嘅生日}，啲仔女就嚟同佢拜壽。

%%K t04-c2-02-1 p
2. --- 細細個嘅彼得仲好細，佢先至三歲。佢見到隻\VUgras{貓} \textsuperscript{(1)} 行
埋嚟。佢想摸佢滑溜溜嘅\VUgras{毛}同條長\VUgras{尾}；佢冇諗到喺嗰啲靚\VUgras{爪仔}
下面收埋咗尖尖嘅惡爪，隨時會抓損佢。

%%K t04-c2-03-1 p
3. --- 拖住佢嘅家姐加比拉認真好多；至於馬素，著住件大\VUgras{褸}
\textsuperscript{(4)} 同條皮\VUgras{腰帶} \textsuperscript{(5)}，就有啲大人款，雖然
條短褲露出佢對\VUgras{襪} \textsuperscript{(6)}。

%%K t04-c2-04-1 p
4. --- 佢哋老豆著咗件厚厚嘅冬天\VUgras{大褸} \textsuperscript{(7)}，有
\VUgras{毛皮領} \textsuperscript{(8)}，\VUgras{袋} \textsuperscript{(9)} 入面放住條
頸巾。佢老婆戴住頂插住\VUgras{羽毛} \textsuperscript{(10)} 嘅氈帽，著住件
\VUgras{短外套} \textsuperscript{(11)}，圍住條\VUgras{毛皮圍巾} \textsuperscript{(12)}，
拎住個\VUgras{暖手筒} \textsuperscript{(13)}。

%%K t04-c2-05-1 p
5. --- 天氣好凍，因為而家係冬天。\VUgras{雪} \textsuperscript{(19)} 落咗成日，啲
屋頂白晒。跟住\VUgras{夜} 一到，就更加凍。天上面，\VUgras{月} \textsuperscript{(17)}
同啲\VUgras{星} \textsuperscript{(18)} 照住，穿透\VUgras{黑暗}。

%%K t04-tit-8 sub
\VUsaut{4.2mm}
\VUpk{10.2pt}{40}{病。}
\VUsaut{3.0mm}

%%K t04-c2-06-1 p
6. --- 嗰個可憐嘅\VUgras{盲人} \textsuperscript{(20)} 真係好唔幸，佢喺\VUgras{藥房}
\textsuperscript{(16)} 前面過緊廣場，由自己隻\VUgras{貴婦狗} \textsuperscript{(21)}
帶路。
\VUgras{女傭} \textsuperscript{(14)} 玉婷由廚房捧上嚟一\VUgras{碗} \textsuperscript{(15)}
放咗好多糖、滾熱辣嘅\VUgras{草藥茶}。

%%K t04-c2-07-1 p
7. --- \VUgras{嫲嫲} \textsuperscript{(23)} 想去檯上面攞自己副\VUgras{眼鏡}
\textsuperscript{(26)} 嚟睇張\VUgras{藥方} \textsuperscript{(27)}——嗰張藥方係
\VUgras{醫生} \textsuperscript{(25)} 啱啱寫嘅——佢就撐住自己支\VUgras{枴杖}
\textsuperscript{(24)} 由張大椅企起身。

%%K t04-c2-08-1 p
8. --- 佢成日\VUgras{唔舒服}、\VUgras{頭暈}、\VUgras{傷風}同\VUgras{牙痛}；佢對
腳好弱，對眼亦都唔靈光。雖然係咁，佢仲好中意笑自己嗰個老\VUgras{醫生}，話佢成日
淨係識開一劑\VUgras{藥水} \textsuperscript{(28)}。今日佢好開心，因為身邊有成家後生。

%%K t04-c2-09-1 p
9. --- 間房閂到密密實實，好舒服。大理石\VUgras{壁爐} \textsuperscript{(29)} 入面燒住
一爐\VUgras{靚火} \textsuperscript{(30)}，喺啱啱點著嘅\VUgras{燈} \textsuperscript{(31)}
柔和嘅\VUgras{光}下面，人人都覺得幸福。

%%K t04-tit-9 sub
\VUsaut{4.2mm}
\VUpk{10.2pt}{40}{阿爺。}
\VUsaut{3.0mm}

%%K t04-c2-10-1 p
10. --- 連\VUgras{阿爺} \textsuperscript{(33)} 自己都唔記得咗自己病緊，唔記得咗
\VUgras{風濕}將佢綁死喺張\VUgras{大椅} \textsuperscript{(34)} 度，仲唔記得咗噖日佢
先至\VUgras{發燒}同\VUgras{暈}過。佢已經唔記得舊年冬天佢生過\VUgras{肺炎}，又畀
一啲沙聲又辛苦嘅\VUgras{咳}折磨到好慘。

%%K t04-c2-10-2 p
但係當時佢真係好難先至好返。之後佢好返少少。不過佢隻腳都仲會痛，佢將佢架喺一個
\VUgras{墊} \textsuperscript{(38)} 上面，個墊就擺喺一張細\VUgras{凳仔}
\textsuperscript{(39)} 度。

%%K t04-c2-11-1 p
11. --- 當佢畀人小心咁裹喺自己件暖笠笠嘅\VUgras{睡袍} \textsuperscript{(35)} 入面
——嗰件睡袍釘住大粒螺鈿\VUgras{鈕} \textsuperscript{(36)}——又有個細細個嘅
\VUgras{孤女} \textsuperscript{(40)} 喺身邊讀報紙畀佢聽，佢戴住白\VUgras{布帽}，
著住藍\VUgras{裙} \textsuperscript{(42)} 同\VUgras{短披肩} \textsuperscript{(41)}
——嗰陣佢就唔會出聲嗌苦。

%%K t04-c2-12-1 p
12. --- 每年，當\VUgras{月曆} \textsuperscript{(43)} 又指住\VUgras{十二月}初一
呢個\VUgras{日子}嗰陣，佢就急住等啲\VUgras{孫}嚟，因為佢知佢哋唔會忘記呢個重要嘅
\VUgras{日子}。

%%K t04-c2-13-1 p
13. --- 佢好感觸咁記得好耐好耐之前，佢親自去同帝國嗰位威名顯赫嘅\VUgras{元帥}
拜壽——嗰位元帥嘅\VUgras{畫像} \textsuperscript{(44)} 而家掛喺牆上面，佢就係佢啲
仔女嘅\VUgras{太公}。
\VUsaut{6.0mm}
\VUfilet{22mm}
